%created on August 18, 2005
%updated on december 19 2005


\chapter {Problem Definition and Formulation}
\label{problem}
In this section, we present two formulations for a deterministic version 
of the LRS problem. In Section \ref{edge-based-form}, we present an edge-based 
formulation of the problem. In Section \ref{set-part-form}, we present a set 
partitioning-based formulation of the problem.

Given a set of candidate facility locations and a set of customer locations, 
the objective of the LRS problem is to select a subset of the facilities, 
to construct a set of delivery routes and to assign routes to vehicles 
in such a way as to minimize total cost subject to the following 
constraints: (i) each customer must be visited exactly once, (ii) each 
route must start at a facility and return to the same facility, (iii) the 
total demand of the customers assigned to a route must be less than the vehicle 
capacity, (iv) the total working hours of a vehicle must be less than the working 
hour limit and (v) the total demand of the customers assigned to a facility 
must not exceed the capacity of the facility.

Throughout this section, we assume that the demand information, distance matrix 
and service times are known and not dynamic. We consider the problem as a delivery 
problem, but the structure of the problem can easily be changed to a pickup problem. 
We assume that each facility has enough vehicles to serve the allocated customers, 
and we consider one class of vehicles, all having the same characteristics and 
working under the same conditions. 

\section{Edge-Based Formulation}
\label{edge-based-form}
To formulate the LRS problem, we define the following sets, parameters
and decision variables.

{\bf Sets}
\begin{eqnarray*}
N & = & \mbox{set of demand nodes} \\
M & = & \mbox{set of candidate facility locations} \\
I & = & \mbox{set of all nodes, $M \cup N$} \\
H & = & \mbox{set of available vehicles} 
\end{eqnarray*}

{\bf Parameters}
\begin{eqnarray*}
FC_j & = & \mbox{daily fixed cost of opening facility $j$, $\forall j \in M$} \\
OC & = & \mbox{hourly operating cost of a vehicle} \\
VF & = & \mbox{total daily fixed cost of a vehicle (including the driver cost)} \\
d_i & = & \mbox{demand of node $i$, $\forall i \in N$} \\
CapF_j & = & \mbox{capacity of facility $j$, $\forall j \in M$} \\
CapV & = & \mbox{capacity of a vehicle} \\
t_{ij} & = & \mbox{travel time between locations $i$ and $j$, $\forall i,j \in I$} \\
s_{i} & = & \mbox{service time at location $i$, $\forall i \in N$} \\
TL & = & \mbox{working hour limit for vehicle and driver}
\end{eqnarray*}

{\bf Decision Variables}
\begin{eqnarray*}
X_{ikh} & = & \left\{ \begin{array}{ll}
                 1  & \mbox{if vehicle $h$ travels directly from location $i$ to location $k$, $\forall i \in I, k \in I, h \in H$}\\
                 0  & \mbox{otherwise}                                                  
                \end{array} 
              \right. \\
Y_{ikh} & = & \mbox{flow from location $i$ to location $k$ carried by vehicle $h$, $\forall i \in I, k \in I, h \in H$} \\
T_{j} & = & \left\{ \begin{array}{ll}
                 1  & \mbox{if facility $j$ is selected, $\forall j \in M$}\\
                 0  & \mbox{otherwise}                                                  
                \end{array}  
              \right. \\
\mbox{Dri}_{h} & = &\left \{ \begin{array}{ll}
                 1  & \mbox{if vehicle $h$ is used, $\forall h \in H$}\\
                 0  & \mbox{otherwise}                                                  
                \end{array}  \right. 
%\mbox{NumH} & = & \mbox{total number of vehicles used} \\
%\mbox{Tot} & = & \mbox{total vehicle hours used}
\end{eqnarray*}

{\bf Formulation}

%We propose the following edge-based formulation for the LRS problem.

\begin{optprog}
{\bf (EB-LRS)} Min & \objective{\sum_{j\in M} FC_j \cdot T_{j}  + VF \cdot \sum_{h \in H} \mbox{Dri}_h  + OC \cdot \sum_{h \in H}\left[\sum_{i \in I}\sum_{k \in I}\left(t_{ik}+s_{k}\right)X_{ikh}\right] } \label{obj-eb} \\
s.t. & \sum_{h\in H}\sum_{k \in I} X_{ikh} & = & 1 &  \forall  i \in N \label{set-part}\\
           & \sum_{k\in I} X_{ikh}-\sum_{k \in I} X_{kih} & = & 0 & \forall i \in I, h \in H \label{enterleave}\\
           & \sum_{h \in H}\sum_{k \in N} Y_{jkh} - CapF_j\cdot T_{j} & \leq & 0 & 
	   \forall j \in M \label{FacCap}\\
           & Y_{ikh} - CapV \cdot X_{ikh} & \leq & 0 & \forall i,k \in I, h \in H \label{VehCap}\\
           & \sum_{h \in H}\sum_{k \in I} Y_{ikh} - \sum_{h \in H}\sum_{k \in I} Y_{kih} & = & - d_i & \forall i \in N \label{flow-conv}\\
           & \sum_{i\in I}\sum_{k \in I}\left(t_{ik}+s_{k}\right)X_{ikh} - TL \cdot \mbox{Dri}_h & \leq & 0 & \forall h \in H \label{timeLim}\\
          % & \mbox{NumH} - \sum_{h \in H} \mbox{Dri}_h & = & 0 & \\
           %& \mbox{Tot} - \sum_{h \in H}\left[\sum_{i \in I}\sum_{k \in I}\left(t_{ik}+s_{k}\right)X_{ikh}\right] / 60 & = & 0 & \\
           & X_{jih} - T_{j} & \leq & 0 & \forall j \in M, i \in I, h \in H \label{Rel-X-T-1}\\
           & X_{ijh} - T_{j} & \leq & 0 & \forall j \in M ,i \in I, h \in H \label{Rel-X-T-2}\\
           & X_{ikh} & \in & \{0,1\} & \forall i,k \in I, h \in H \label{x-var}\\
           & T_{j}  & \in & \{0,1\} & \forall j \in M \label{t-var}\\
           & Dri_{h} & \in & \{0,1\} &  \forall h \in H \label{dri-var}\\
           & Y_{ikh} & \geq &  0 & \forall i,k \in I, h \in H \label{y-var} 
         %  & \mbox{Tot} & \geq & 0 
\end{optprog}

%think about the edge based formulation arcs from facilities to other..

The objective function (\ref{obj-eb}) seeks to minimize the total cost, which  
includes the fixed cost of the selected facilities, the fixed cost of the 
drivers and vehicles and the operating cost of the vehicles. Constraints 
(\ref{set-part}) specify that exactly one vehicle must travel from demand node $i$
to some other node. Constraints (\ref{enterleave}) say that each vehicle should enter 
and leave a node an equal number of times. For demand nodes, constraints 
(\ref{set-part}) and (\ref{enterleave}) ensure that if a vehicle enters a 
demand node, eventually it leaves this node and therefore that each demand 
node is served exactly once. For facility nodes, the number of visits by any 
vehicle can be more than one, since a vehicle can cover multiple routes. 
Constraints (\ref{FacCap}) are capacity constraints that ensure that the 
capacity of a facility is not exceeded by flows to demand nodes. 
Constraints (\ref{VehCap}) are vehicle capacity constraints and define the 
relationship between the binary variable $X_{ikh}$ and flow variable 
$Y_{ikh}$. Each constraint ensures that any flow between any node pair 
cannot be greater than the load capacity of a vehicle. Constraints (\ref{flow-conv})
require conservation of flow at each demand node based on the assumption that the 
problem is a delivery problem. This constraint should be changed if the type of
service is pickup instead of delivery. Constraints (\ref{VehCap}) and (\ref{flow-conv}) 
together ensure that the vehicle capacity is not violated in any route of the vehicle.
In combination with constraints (\ref{set-part}) and (\ref{enterleave}), 
constraints (\ref{flow-conv}) form the travelling salesman tours, as in 
\citet {Nambiar}.  Constraints (\ref{timeLim}) limit the 
total working time of a vehicle to the working hour limit. 
%Constraints (8) and (9) simply calculate the total number of used vehicles and the vehicle 
%hours, respectively. 
Constraints (\ref{Rel-X-T-1}) and (\ref{Rel-X-T-2}) link variables $X_{ijh}$ 
and $T_{j}$; facility $j$ must be selected if vehicles enter or leave
from it. Constraints (\ref{x-var}), (\ref{t-var}), (\ref{dri-var}) and (\ref{y-var}) 
are the set of integrality and non-negativity requirements on the variables.    

%%%%%%%%%%%%%%%%%%%%%%%%%%%%%%%%%%%%%%%%%%%%%%%%%%%%%%%%%%%%%%%%%%
%%%%%%%%%%%%%%%%%%%%%%%%%%%%%%%%%%%%%%%%%%%%%%%%%%%%%%%%%%%%%%%%%
%PROVE NP-HARD 
To show that the LRS problem is NP-hard, we start from the fact that the
multi-depot vehicle routing problem (MDVRP) is NP-hard (\citet{Bodin2,Lenstra}). 
The MDVRP aims to find a minimum-cost set of vehicle routes, each one starting from a
depot and terminating at the same depot. Each customer must be served exactly once. 
There are specified lower and upper bounds for the number of vehicles at each depot
and vehicles have capacities. To show that the LRS problem is NP-hard, we show that the 
MDVRP is a special case of the LRS problem.  We will consider an instance of the LRS 
problem with $N$ customer nodes and $M$ facilities and at most $N$ vehicles at 
each depot. In the created instance, to eliminate the location part of the LRS problem, 
the fixed cost of opening a facility is set to zero. Since the MDVRP considers only the 
operating costs of vehicles, the fixed cost of a vehicle (in the LRS problem) is assumed 
to be zero for the instance. 
%To be able to serve all customers allocated to a facility, 
The total working hour limit of a vehicle is 
eliminated, or equivalently the total working hour limit of the vehicle is set to infinity.
The created instance of the  LRS problem is equivalent to an instance of the MDVRP 
with $N$ customer nodes and $M$ facilities, each of which has $0$ to $N$ vehicles.
The solution of the special case of the LRS problem can easily be translated to
a solution for the MDVRP. In the solution of the MDVRP, the allocation of the 
customers to the facilities, the vehicle routes and the total cost are the same as in 
the solution of the LRS instance. 
%The only difference is the number of vehicles used in the solution. 
%In the MDVRP, each route is assigned to one vehicle, therefore
%the number of vehicles used is equal to the number of routes. 
%However, in the created instance of the LRS problem, there is one vehicle at each facility and 
%all the routes originated from a facility are assigned to the vehicle located at 
%the facility. 
Since we have shown that the MDVRP is a special case of the LRS 
problem and the MDVRP is NP-hard, we can conclude that the LRS problem is NP-hard. 
 
%If both problems have the same distance matrix and same operating cost, then solving the LRS 
%problem yields a solution for the MDVRP.
 
%Each facility has a capacity of value $FCap$ and has 1 vehicle with capacity $VCap$ and no total working hour 
%limit, which means that we set the 
%total working hour limit of each vehicle to infinity.  This instance for 
%LRS problem can be converted to an instance of MDVRP with M facilities 
%and N customer nodes in which each facility has 0 to 
%N vehicles with capacity $VCap$. If both problems have the same 
%distance matrix and same operating cost, then solving the LRS problem 
%is equavalent to solving the MDVRP. 

%We define a special case of LRS problem,
%solution of which is equivalent to solution of MDVRP.
%%%%%%%%%%%%%%%%%%%%%%%%%%%%%%%%%%%%%%%%%%%%%%%%%%%%%%%%%%%%%%%%%%%
%[Dr. Berger deleted this part] ??
%The allocation of customers
%to facilities, vehicle routes and the total costs are same. 
%The only difference is in MDVRP at most one route is assigned to
%each vehicle although in LRS problem it is possible to assign
%multiple routes to a vehicle. Thus, if multiple vehicles are used to
%serve customers in any depot, while converting the 
%solution of MDVRP to the solution of LRS, 1 vehicle is assigned to these 
%multiple routes in that depot. This assignment is possible since there is no working hour 
%limit for the vehicles. Since a special case of LRS problem is equivalent to
%a MDVRP, LRS problem is NP-Hard.           

%%%%%%%%%%%%%%%%%%%%%%%%%%%%%%%%%%%%%%%%%%%%%%%%%%%%%%%%%%%%%%%%%%%
The EB-LRS formulation is a three-index commodity flow model, and in the worst case, 
it has $O(|N|^4)$ binary vaiables and $O(|N|^4)$ constraints. Therefore,
there are more than 6 million constraints and binary variables for an 
instance with $50$ customers. We could develop a branch-and-bound or
a branch-and-cut algorithm for the EB-LRS problem, however, we suspect
that such an approach will not yield good results, given the difficulty 
others have had in solving related problems with similar commodity formulations. 
For example, exact algorithms for the LRP generally are developed 
for two-index formulations of the problem, and none of the commodity 
flow models developed for the VRP is solved to optimality in the
literature. Therefore, we instead are propose a set 
partitioning-based approach, which has been successful for other hard 
combinatorial problems. 
 

%Besides the fact that the edge-based formulation, EB-LRS has many binary variables 
%and many constraints, the EB-LRS is a three-index commodity flow model. 
%Exact algorithms for the LRP are generally developed for two-index formulations of the 
%problem.  Also, none of the commodity flow models developed for the VRP is solved to optimality in the
%literature. Furthermore, if there are side constraints for the vehicle routes such as 
%vehicle capacity, route length restriction and time windows, the decomposition algorithms 
%are common in the literature. Based on these observations, we develop an equivalent set 
%partitioning formulation for the LRS problem and we solve the set partitioning-based 
%formulation using a branch-and-price algorithm. 

%[moved to some other place]Although branch-and-price has been applied successfully to some of the 
%large scale optimization problems such as crew scheduling, machine
%scheduling, facility location and routing, since the mid-1990s, there are 
%many difficulties associated with the algorithm and the success of the 
%algorithm highly depends on the structure of the problem. \citet{Barnhart}
%provide an overview of the general methodology and many of the challenges.

%the formulation is not easily solved using conventional MIP methods


%%%%%%%%%%%%%%%%%%%%%%%%%%%%%%%%%%%%%%%%%%%%%%%%%%%%%%%%%%%%%%%%%%%%%%%%

\section{Set Partitioning Formulation}
\label{set-part-form}
\subsection{Pairing Formulation}
In this section, we introduce a set partitioning-based formulation of 
the LRS problem. The formulation uses the idea of a {\em pairing}, which is 
adapted from the pairing concept in the crew scheduling literature 
(e.g., \citet{Vance}, \citet{Anbil} and \citet{Desrosiers}). In the 
crew scheduling literature, a pairing is a sequence of duty periods and represents  
a possible schedule for a single crew. In the LRS context, a pairing is a 
set of routes that can be served sequentially by one vehicle within the 
vehicle's working hour limit. Figure \ref {fig:exppair} shows two possible 
pairing constructions. Each pairing represents the schedule of one vehicle. 
One of the pairings includes three paths serving eight customers in total 
and the other pairing has two paths with three customers in total. We assume that
the cost of a pairing is the sum of the costs of the component routes plus a
fixed cost for the vehicle. Since the cost of a route depends on the sequence 
of nodes visited, the pairing information implicitly includes sequencing information
for each route. However, when there are no time windows, the ordering 
of the routes within a pairing is arbitrary. 

\begin{figure}[h]
  \centering
  \includegraphics[scale=0.5]{figs/pairexp2.eps}
  %\fbox{ \includegraphics[scale=0.5]{figure/pair2.eps}}
  \caption{Example of pairings} \label{fig:exppair}
\end{figure}

A pairing is feasible if (i) total demand of customers in each route included in 
the pairing is at most the vehicle capacity, (ii) total travel time of
the pairing is at most the vehicle working hour limit, and (iii) each
customer included in the pairing is visited once.  Given a set of candidate 
facility locations and the set of feasible pairings that can be constructed 
for each candidate facility, the objective of the LRS problem is to select a 
subset of the facilities and an associated subset of pairings in such a way 
as to minimize the total cost subject to the following constraints: (i) each 
customer must be covered by exactly one pairing, (ii) the total demand of the 
customers assigned to a facility must not exceed the capacity of the facility, 
and (iii) vehicles at a facility can be used only if the facility is open. 
%(iii) the total time of the routes in the pairing 
%should not exceed the daily time limit. 

As in Section~\ref{edge-based-form}, we let $N$ denote the 
set of demand nodes, $M$ denote the set of candidate facilities and 
$I$ denote the set of all nodes. Then, we define the following
additional notation. Associated with each candidate
facility $j \in M$ is a set $P_j$, which is the set of all feasible 
pairings in which the vehicle starts and ends at facility $j$.  

{\bf Parameters}
\begin{eqnarray*}
a_{ipj} & = & \left \{  \begin{array}{ll}
                 1 & \mbox{if demand node $i$ is in pairing $p$ of facility $j$, $\forall i \in N, j \in M, p \in P_{j}$ } \\
                 0 & \mbox{otherwise}
                         \end{array}
              \right. \\
C_{jp} & = & \mbox{cost of pairing $p$ associated with facility $j$, $\forall j \in M, p \in P_{j}$}
\end{eqnarray*}
%$nREQ$ is simply calculated by using the
%demands of the customers and the capacity of the facilities. Assuming that all
%facilities have same capacity value, $CapF$, then 
%\begin{eqnarray}
%nREQ =\displaystyle \left \lceil \frac{\sum_{i \in N }d_i}{CapF} \right \rceil \nonumber 
%\end{eqnarray}   
%If the capacity is different for each facility, then $nREQ$ is calculated
%by assigning the total demand to the facilities that are sorted 
%in decreasing order of their capacities. 
%    & \sum_{j \in M} T_{j} & \geq & nREQ \\
%  & Z_{jp} & \leq & T_{j}  & \forall j \in M, \forall p \in P_{j} \\

{\bf Decision Variables}
\begin{eqnarray*}
Z_{jp} & = & \left \{ \begin{array}{ll}
                 1  & \mbox{if pairing p is selected for facility j, $\forall p \in P_{j}$ and  $j \in M$ }\\
                 0  & \mbox{otherwise}                                                  
                \end{array}
       \right. \\
T_{j} &  = & \left \{ \begin{array}{ll}
                 1  & \mbox{if facility j is selected, $\forall j \in M$ }\\
                 0  & \mbox{otherwise}                                                  
                \end{array}
       \right.
\end{eqnarray*}

{\bf Formulation}
\begin{optprog}
{\bf(SPP-LRS)} Minimize & \objective{\sum_{j \in M} FC_j T_{j} + \sum_{j \in M} \sum_{p\in P_{j}} C_{jp} Z_{jp}} \label{spp-obj}\\
subject to & \sum_{j \in M} \sum_{p \in P_{j}} a_{ipj} Z_{jp} & = & 1 & \forall i \in N \label{spp_con1} \\
           & \sum_{p \in P_j} \sum_{i \in N} a_{ipj} d_i Z_{jp} & \leq & CapF_{j} T_{j} & \forall j \in M \label{spp_con2} \\
           & Z_{jp} &  \in & \{0,1\} & \forall j \in M, \forall p \in P_{j} \label{binary1} \\
           & T_{j}  & \in & \{0,1\} & \forall  j \in M \label{binary2}
\end{optprog} 
 
The objective function (\ref{spp-obj}) seeks to minimize the total cost, which includes the fixed 
cost of the selected facilities and the cost of the selected pairings. We note that the cost $C_{jp}$ 
of a pairing includes both the vehicle fixed costs and variable costs. Constraints (\ref{spp_con1}) 
guarantee that each demand location is served by exactly one pairing. Constraints (\ref{spp_con2}) 
are capacity constraints that ensure that the capacity of a facility is not exceeded by the selected 
pairings. Constraints (\ref{binary1}) and (\ref{binary2}) are standard binary restrictions on the 
variables. The structure of the SPP-LRS is very similar to the CFLP with single sourcing. In the
CFLP with single sourcing, only pairings with one customer are considered and assigned to capacitated
facilities. However, in the LRS problem, pairings with multiple customers must be taken into account. 

%While customers are assigned to capacitated facilites in the CFLP with single sourcing, pairings are
%assigned to capacitated facilites in the SPP-LRS






\subsection {Dantzig Wolfe Decompostion of EB-LRS} 
The set partitioning formulation (SPP-LRS) 
can be obtained by applying Dantzig Wolfe decomposition to the edge-based formulation of the LRS problem 
(EB-LRS). We can define problem (EB-LRS) as 
\begin{optprog}
{\bf(EB-LRS')} Minimize & \objective{CX'} \nonumber\\
subject to & AX' & \leq & D_1 \label{sys1}\\
           & BX' & \leq & D_2 \label{sys2}
\end{optprog} 

\noindent where 
\begin{itemize}
\item system (\ref{sys1}) represents the system of constraints, \{Const.(\ref{set-part}), Const.(\ref{FacCap}), 
  Const.(\ref{Rel-X-T-1}), Const.(\ref{Rel-X-T-2}), Const.(\ref{x-var}), Const.(\ref{t-var}), Const.(\ref{y-var})\},
\item system (\ref{sys2}) represents the set, \{ Const.(\ref{enterleave}),
  Const.(\ref{VehCap}), Const.(\ref{flow-conv}), Const.(\ref{timeLim}), Const.(\ref{x-var}), Const.(\ref{dri-var}) 
  and Const.(\ref{y-var})\},
\item $X'=(X,Y,T)$.
\end{itemize}

%(\ref{sys1}) represents the system of constraints, \{Const.(2), Const.(4), 
%Const.(10), Const.(11), Const.(12), Const.(13), Const.(15)\}, 
%system (\ref{sys2}) represents the set, \{ $\leq$ part of Const.(2), Const.(3),
%Const.(5), Const.(6), Const.(7), Const.(8), Const.(9), Const.(12), Const.(14), 
%Const.(15) and Const.(16)\} and $X'=(X,Y,T)$. 

The set of feasible solutions for the system of constraints in $BX'  \leq  D_2$ 
(system (\ref{sys2})) is equal to the set of solutions of the LRS problem, which are 
feasible with respect to the vehicle capacity constraints and the vehicle time 
limit constraints. The system  $AX'  \leq  D_1$ (system (\ref{sys1})) selects
the solution that is feasible with respect to the facility capacity constraints
and set partitioning constraints (\ref{set-part}) among the solutions provided by  
system (\ref{sys2}). 

Instead of system (\ref{sys2}), we can define system $B'X'  \leq  D_2'$, which
includes the system $BX'\leq  D_2$, e.g. $\{X'|BX'  \leq  D_2\}\subseteq \{X'|B'X'  \leq  D_2'\}$.  
\begin{optprog}
 $B'X' \leq D_2'= \{$  & \sum_{h \in H_j} \sum_{k \in I} X_{ikh} & \leq & 1 &  \forall  i \in N , j \in M \label{rel-spp}\\
           & \sum_{k\in I} X_{ikh}-\sum_{k \in I} X_{kih} & = & 0 & \forall i \in I, j \in M, h \in H_j \nonumber\\
           & Y_{ikh} - CapV \cdot X_{ikh} & \leq & 0 & \forall i,k \in I, j \in M, h \in H_j\nonumber \\
           & \sum_{k \in I} Y_{ikh} - \sum_{k \in I} Y_{kih} + d_i \sum_{h \in H_j} \sum_{k \in I}X_{ikh} & = & 0 & \forall i \in N, j \in M \label{flowvar} \\
           & \sum_{i \in I}\sum_{k \in I}\left(t_{ik}+s_{k}\right)X_{ikh} - TL\cdot Dri_h  & \leq & 0 & \forall j \in M, h \in H_j\nonumber\\
           & X_{ikh} & \in & \{0,1\} & \forall i,k \in I, j \in M, h \in H_j\nonumber\\
           & Y_{ikh} & \geq &  0 & \forall i,k \in I, j \in M, h \in H_j \nonumber \\
           & Dri_h  & \in & \{0,1\} & \forall j \in M, h \in H_j\nonumber\\
\} \nonumber
\end{optprog}

Where $H_j$ is the set of vehicles located at facility $j$. The constraints  
in the system $B'X'  \leq  D_2'$ define the feasible pairings for each facility
with respect to the vehicle capacity constraints and the vehicle time limit constraints. 
The constraints (\ref{rel-spp}) are relaxed forms of our set partitioning 
constraints (\ref{set-part}). We need the relaxed form because a feasible 
pairing does not need to visit all of the customer nodes but it may not visit a 
customer node more than once. Because of the same reason, in constraint 
(\ref{flowvar}), we multiply $d_i$ with $\sum_{h \in H_j}\sum_{k \in I}X_{ikh}$. If the 
pairing does not visit the customer node $i$, then  the flow into the node $i$ 
minus the flow out of the node $i$ must be zero. Otherwise, the summation 
$\sum_{h \in H_j}\sum_{k \in I}X_{ikh}$ is one and the difference between 
the flow into the customer and out of the customer must equal to the demand of 
the customer.


Since each constraint in $B'X'  \leq  D_2'$ is defined for each facility $j \in M$, we can decompose 
the system $B'X'  \leq  D_2'$ into $M$ independent subproblems 
$B'_jX_j' \leq D_{2j}'$,  one for each facility $j \in M$. Each subproblem 
generates feasible pairings for a facility. Similar to our SPP-LRS formulation, we define  
$P_j$ as the set of feasible pairings at facility $j$ and $a_{pj}$ as the column 
vector such that  
\begin{eqnarray*}
a_{ipj} & = & \left \{  \begin{array}{ll}
                 1 & \mbox{if node $i$ is in pairing $p$ of facility $j$, $\forall i \in N, j \in M, p \in P_{j}$ } \\
                 0 & \mbox{otherwise}
                         \end{array}
              \right. \\
C_{jp} &=&\mbox{cost of the pairing pairing $p$  of facility $j$ (column $a_{pj}$), $\forall j \in M, p \in P_{j}$} 
\end{eqnarray*} 
Since we are solving one subproblem for each facility, we define our node set 
as $I^j=N \cup \{j\}$. We will create one feasible pairing, so we can drop the 
index h from the original variable definition and define variables 
$(X_{ik},Y_{ik})^{N+1,N+1}$:
\begin{eqnarray*}
X_{ik} & = & \left\{ \begin{array}{ll}
                 1  & \mbox{if arc (i,k) is selected $\forall i \in I^j, k \in I^j$}\\
                 0  & \mbox{otherwise}                                                  
                \end{array} 
              \right. \\
Y_{ik} & = & \mbox{flow from location $i$ to location $k$, $\forall i \in I^j, k \in I^j$}
\end{eqnarray*}
Each $(X_{ik},Y_{ik})^{N+1,N+1}$ vector represents one feasible pairing for facility $j$.
We can define the set $P_j$ by using the system $B'_jX_j' \leq D_{2j}'$:

\begin{optprog}
$P_j= \{(X_{ik},Y_{ik})^{N+1,N+1}:$  & \sum_{k \in I^j} X_{ik} & \leq & 1 &  \forall  i \in N \nonumber \\
           & \sum_{k\in I^j} X_{ik}-\sum_{k \in I^j} X_{ki} & = & 0 & \forall i \in I^j \nonumber\\
           & Y_{ik} - CapV \cdot X_{ik} & \leq & 0 & \forall i,k \in I^j \nonumber \\
           & \sum_{k \in I^j} Y_{ik} - \sum_{k \in I^j} Y_{ki} + d_i \sum_{k \in I^j}X_{ik} & = & 0 & \forall i \in N \label{f:con} \\
           & \sum_{i,k \in I^j}\left(t_{ik}+s_{k}\right)X_{ik} - TL  & \leq & 0 \nonumber\\
           & X_{ik} & \in & \{0,1\} & \forall i,k \in I^j \nonumber\\
           & Y_{ik} & \geq &  0 & \forall i,k \in I^j \nonumber
\}
\end{optprog}
%The first constraint in the definition of $P_j$ is the relaxed form of our set 
%partititoning constraint (\ref{set-part}). We have to use the relaxed form because in a 
%feasible pairing it is not necessary to visit all customer nodes, but still a 
%customer node can be visited at most once. 
 
Define 
\begin{eqnarray*}
 \lambda^j_p & = & \mbox{the number of times pairing $p$ associated with facility $j$ (column $a_{pj}$) is}\\ 
             & & \mbox{selected in the solution of (EB-LRS'),} \ \ \forall j \in M, p \in P_j
\end{eqnarray*}
To construct the Dantzig Wolfe Decomposition of the problem (EB-LRS'), we write $(X,Y)^{pj}$, 
the vector of variables defining pairing $p$ for facility $j$, in terms of column $a_{pj}$ and
variable $\lambda^j_p$. Let $X^{B'X'\leq D_2'}$ define all feasible points satisfying 
the system $B'X'\leq D_2'$. Then we have the following relations: 

%Old version
%\begin{optprog}
%$X^{BX\leq D_2}$ &= \{(X,Y)^j_{j \in M} : (X,Y)^j &\in & P_j \ \ \forall j \in M \} \nonumber\\
%                 &= \{(X,Y)^j_{j \in M} : X^j_{ik} &\leq & a_{ipj}\lambda^j_p \leq \lambda^j_p & \forall  i,k \in I^j, j\in M, p \in P_j \nonumber \\
%                 & \sum_{k \in I^j} X^j_{ik} &=& a_{ipj}\lambda^j_p  & \forall  i \in I^j, j\in M, p \in P_j \nonumber \\
%                  & -\sum_{k \in I^j} Y^j_{ki} &=& -\sum_{i \in N}d_i a_{ipj}\lambda^j_p  & \forall j\in M, p \in P_j  \label{f:con:tran}\\ 
%                 & \sum_{p \in P_j}\lambda^j_p &\leq & |H^j|  & \forall j \in M \label{conv}\\
%                 & \lambda^j_p &\in & \{0,1\} & \forall j \in M, p \in P_j \}\nonumber  
%\end{optprog}

[{\bf The following part was updated}]

\begin{optprog}
$X^{B'X\leq D_2'}$ &= \{(X,Y)^j_{j \in M} : (X,Y)^j &\in & P_j \ \ \forall j \in M \} \nonumber\\
                 &= \{(X,Y)^j_{j \in M} : X^j_{jk} &\leq & \sum_{p \in P_j}a_{kpj}\lambda^j_p  & \forall  k \in N, j\in M \label{forXTrel1}\\
                 & X^j_{kj} &\leq & \sum_{p \in P_j}a_{kpj}\lambda^j_p  & \forall  k \in N, j\in M \label{forXTrel2}\\
                 & \sum_{k \in I^j} X^j_{ik} &=& \sum_{p \in P_j}a_{ipj}\lambda^j_p  & \forall  i \in N, j\in M  \label{forsetpart} \\
                  & -\sum_{i \in N} Y^j_{ji} &=& -\sum_{i \in N}\sum_{p \in P_j}d_i a_{ipj}\lambda^j_p  & \forall j\in M  \label{f:con:tran}\\ 
                 & \sum_{p \in P_j}\lambda^j_p &\leq & |H^j|  & \forall j \in M \label{conv}\\
                 & \lambda^j_p &\in & \{0,1\} & \forall j \in M, p \in P_j \}\nonumber  
\end{optprog}

[{\bf Explanation:}

-We could not give the equivalence of variables $X^j_{ik}$ and $Y^j_{ik}$ (e.g., $X^j_{ik}=?$ and $Y^j_{ik}=?$). But to represent 
system (\ref{sys1}) in terms of variables $\lambda$, we need relations (\ref{forXTrel1}), (\ref{forXTrel2}), (\ref{forsetpart})
and (\ref{f:con:tran}). 

-Relations (\ref{forXTrel1}) and (\ref{forXTrel2}): We need them to represent Const.(\ref{Rel-X-T-1}), Const.(\ref{Rel-X-T-2}). 
They are true, since if $X^j_{jk}=0$ then $\sum_{p \in P_j}a_{kpj}\lambda^j_p$ may be 0 or 1
and if $X^j_{jk}=1$ then $\sum_{p \in P_j}a_{kpj}\lambda^j_p$ will be 1. Also, to write constraints (\ref{Rel-X-T-1}) and (\ref{Rel-X-T-2}), 
we need $\sum_{p \in P_j}a_{kpj}\lambda^j_p \leq T_j$.          
Based on this new relation, the constraint (\ref{rel-1}) should be updated to
$\sum_{p \in P_j}a_{ipj}\lambda^j_p  \leq  T_{j} \forall j \in M, \forall i\in N $

-Relations (\ref{forsetpart}): We need them to represent Const.(\ref{set-part}). 

-Relations (\ref{f:con:tran}): We need them to represent Const.(\ref{FacCap}).

\begin{figure}[h]
  \centering
  \includegraphics[scale=1]{figs/danzigfig.eps}
\end{figure} 

For example based on the figure above, we can write flow conservation equation (\ref{f:con}) for all $i \in N$:
\begin{eqnarray*}
(Y_{13}+Y_{14}+Y_{12}+Y_{15}+Y_{10})-(Y_{01}+Y_{31}+Y_{21}+Y_{51}+Y_{41})+d_1\sum{k \in I^j}X_{1k}=0 \\
(Y_{31}+Y_{34}+Y_{32}+Y_{35}+Y_{30})-(Y_{03}+Y_{13}+Y_{23}+Y_{53}+Y_{43})+d_3\sum{k \in I^j}X_{3k}=0 \\
(Y_{43}+Y_{41}+Y_{42}+Y_{45}+Y_{40})-(Y_{04}+Y_{34}+Y_{24}+Y_{54}+Y_{14})+d_4\sum{k \in I^j}X_{4k}=0 \\
(Y_{23}+Y_{24}+Y_{21}+Y_{25}+Y_{20})-(Y_{02}+Y_{32}+Y_{12}+Y_{52}+Y_{42})+d_2\sum{k \in I^j}X_{2k}=0 \\
(Y_{53}+Y_{54}+Y_{52}+Y_{51}+Y_{50})-(Y_{05}+Y_{35}+Y_{25}+Y_{15}+Y_{45})+d_5\sum{k \in I^j}X_{5k}=0 
\end{eqnarray*}
When we sum equation (\ref{f:con}) for all $i \in N$, we get $\sum_{i \in N}Y_{i0}-\sum_{i \in N}Y_{0i}+\sum_{i \in N}d_i \sum_{k \in I^j}X_{ik}=0$.
We know that $\sum_{i \in N}Y_{i0}=0$. 

]


 
Equality (\ref{f:con:tran}) is obtained by adding the flow conservation equation (\ref{f:con}) for all $i \in N$ and
using the information that the vehicle turns back to depot with zero load because of the characteristics 
of equations (\ref{f:con}) in system (\ref{sys2}).    

By using the transformation of $(X,Y)^{pj}$ to columns $a_{pj}$ and variable $\lambda^j_p$, we can write the
master problem, which is the system (\ref{sys1}), $AX'\leq D_1$, with the objective function. To write the objective 
function, we use the cost definition of the pairing, that is $C_{jp}=VF+OC \sum_{i,k \in I^j}(t_{ik}+s_k)X_{ik} $. 

\begin{optprog}
{\bf (Master-DW)} \nonumber\\
Minimize & \objective{\sum_{j \in M} FC_j T_{j} + \sum_{j \in M} \sum_{p\in P_{j}} C_{jp} \lambda^j_p} \nonumber\\
subject to & \sum_{j \in M} \sum_{p \in P_{j}} a_{ipj} \lambda^j_p & = & 1 & \forall i \in N \label{spar} \\
           & \sum_{p \in P_j} \sum_{i \in N} a_{ipj} d_i \lambda^j_p  & \leq & CapF_{j} T_{j} & \forall j \in M \label{rel-2}  \\
           & \sum_{p \in P_j}a_{ipj}\lambda^j_p & \leq & T_{j}  & \forall j \in M,  \forall i\in N \label{rel-1} \\
           & \sum_{p \in P_j} \lambda^j_p  & \leq & H_j & \forall j \in M \label{conv2} \\
           & \lambda^j_p  & \in & \{0,1\} & \forall  j \in M, \forall p \in P_{j} \label {conv3}\\
           & T_{j}  & \in & \{0,1\} & \forall  j \in M \nonumber
\end{optprog} 

%To write Master-DW in terms of new variables, along with the definition of $X^{B'X'\leq D_2'}$, we
%have used some logical equalities. For example, in constraint (\ref{spar}) we have used the 
%transformation $\sum_{h \in H} \to \sum_{j \in M} \sum_{p \in P_{j}}$ along with the equality 
%$\sum_{k \in I^j} X^j_{ik} = a_{ipj}\lambda^j_p $. 

%%%%%%%%%%%%%%%%%%%%%%%%%%%%%%%%%%%%%%%%%%%%%%%%%%%%%%%%%%%%%%%%%%%%%%%%%%%%%%%%%%%%%%%%%%%%
% rewrite this part    
Constraint (\ref{rel-1}) is implied by Constraint (\ref{rel-2}). But, constraint (\ref{rel-1})
is a valid inequality, thus we can add this constraint that is representing the relationship 
between variable $T_j$ and $\lambda^j_{p}$ to our set partitioning-based formulation. 
In SPP-LRS, that constraint will be 
$\sum_{p \in P_j}a_{ipj}Z_{jp} \leq  T_{j}\ \ \forall j \in M, \forall i \in N$. 
%$Z_{jp} \leq  T_{j}\ \ \forall j \in M, \forall p \in P_{j}$. 

Constraints (\ref{conv2}) and (\ref{conv3}) are convexification constraints. Constraint 
(\ref{conv2}) states the upper bound for the number of vehicles for each facility. 
It is true that we can add an upper bound for the number of vehicles for
each facility to our SPP-LRS formulation. Therefore, Master-DW is same with 
our SPP-LRS formulation with some additional constraints.     

%%%%%%%%%%%%%%%%%%%%%%%%%%%%%%%%%%%%%%%%%%%%%%%%%%%%%%%%%%%%%%%%%%%%%%%%%%%%%%%%%%%%%%%%%
%%%%%%%%%%%%%%%%%%%%%%%%%%%%%%%%%%%%%%%%%%%%%%%%%%%%%%%%%%%%%%%%%%%%%%%%%%%%%%%%%%%%%%%

\subsection{Strengthening the Formulation}
\label{StrongForm}
In this section, we discuss how we can strengthen the SPP-LRS formulation by adding 
some valid inequalities. First, we can introduce a lower bound on the total 
number of facilities selected. Let us define parameter  $nREQ$ as the minimum 
number of required (selected) facilities. For any instance, $nREQ$ can be calculated 
from the demands of the customers and the capacities of the facilities. There are 
two possible cases:

%\begin{itemize}
%\item 
\noindent {\bf Case 1.} If all of the facilities have the same capacity value, $CapF$, then  $nREQ$ 
is calculated by dividing the total demand by the capacity and rounding up:
\begin{equation}
nREQ =\left \lceil \frac{\sum_{i \in N }d_i}{CapF} \right \rceil\nonumber
\end{equation}
%\item 
\noindent {\bf Case 2.} If at least one facility has a different capacity value,  
then a lower bound for the number of facilities can be calculated using a greedy 
algorithm: Sort the facilities in non-increasing order of their capacities. Then
assign the total demand to the facilities in order until all the demand is assigned.  
Let $(j_1,j_2,..,j_M)$ be the order of the facilities such that 
$CapF_{j_1}\geq CapF_{j_2} \geq ...\geq CapF_{j_M}$. Then,
\begin{eqnarray*}
 nREQ =k^*=argmin_{\{k \in 1..M\}} \left (\sum_{t=1}^{k}CapF_{j_t}\geq \mbox{total demand} \right )  
\end{eqnarray*}
%\end{itemize}
After determining the lower bound, $nREQ$, we can add to the SPP-LRS formulation 
the constraint
\begin{eqnarray}
\sum_{j \in M} T_{j} \geq nREQ \label{lb-open-fac}
\end{eqnarray}
Constraint (\ref{lb-open-fac}) is not required to define the LRS problem. However, it is a simple
lower bound for the total number of facilities that should be open, and it might eliminate solutions
in which a fractional number of facilities is used. With the constraint, the LP solution may be closer 
to the IP solution in terms of the $T_j$ variables. The following proposition can be stated:

%In the LP solution of the model,  
%if $\sum_{j \in M} T_{j}$ is fractional, then there is at least one fractional $T_j$ variable.
%In addition, if  $\sum_{j \in M} T_{j}$ is strictly 
%less than $nREQ$, then constraint (\ref{lb-open-fac}) forces a non-zero fractional $T_j$ variable to 
%be $1$. Therefore, with the constraint, the LP solution may be closer to the IP solution in terms of 
%the $T_j$ variables. The following proposition can be stated:


%lower bound for the optimal total cost may improve significantly with 
%constraint (\ref{lb-open-fac}). 
 

\begin {prop} If all of the facilities have the same capacity value, $CapF$, and 
$\left (\frac{\sum_{i \in N }d_i}{CapF}\right)$ is fractional, or facilities have
different capacity values and $\left (\frac{\sum_{i \in N }d_i}{max_{j \in M}\{CapF_j\}}\right)$
is fractional or strictly less than $nREQ$, then the 
Chv\'atal-Gomory rank (CG rank) of constraint (\ref{lb-open-fac}) is 1. 
\end {prop}

\begin{proof}

Recall that the CG rank of any original inequality and any 
inequality dominated by a nonnegative linear combination of the original inequalities 
is 0.  An inequality which can be generated as a CG inequality based on rank 0 inequalities 
but is not rank 0 itself has CG rank 1. 

{\bf Case 1:} The facilities have the same capacity value, $CapF$.

To prove the proposition, we compute the
summation of constraints (\ref{spp_con2}) over all facilities $j \in M$:
\begin{eqnarray}
 \sum_{j \in M} \sum_{p \in P_j} \sum_{i \in N} a_{ipj} d_i Z_{jp} & \leq & CapF \sum_{j \in M} T_{j} \label{sum-spp_con2}
\end{eqnarray} 
Then, if we substitute constraint (\ref{spp_con1}) in (\ref{sum-spp_con2}) and divide by $CapF$, we get:
\begin{eqnarray}
 \sum_{j \in M} T_{j} &\geq &\frac {\sum_{i \in N}d_i}{CapF} \label{sum2-spp_con2}
\end{eqnarray} 

The $T_j$ variables are binary, so $\sum_{j \in M} T_{j}$ must be integer. 
Therefore, if we let 
$nREQ =\left \lceil \frac{\sum_{i \in N }d_i}{CapF} \right \rceil$, then constraint
(\ref{lb-open-fac}) must be valid. If the right-hand side of (\ref{sum2-spp_con2}) 
is fractional, the constraint (\ref{lb-open-fac}) has CG rank 1. Otherwise, its rank is 0,
since it can be defined as the non-negative linear combination of the original inequalities.
 

{\bf Case 2:} The facilities have different capacity values. 

To generalize this statement for the problems in which the facility
capacities might be different in value, we reorder the facilities as
$T_{[1]},T_{[2]},..,T_{[M]}$ such that $CapF_{[1]} \geq CapF_{[2]} \geq 
..\geq CapF_{[M]}$. Then, inequality (\ref{sum-spp_con2}) becomes 
\begin{eqnarray}
\sum_{j \in M} \sum_{p \in P_j} \sum_{i \in N} a_{ipj} d_i Z_{jp} \leq   \sum_{j \in M} CapF_{[j]}T_{[j]} \label{nonhomo-sum}
\end{eqnarray}  
Again, if we substitute in (\ref{nonhomo-sum}) using constraint (\ref{spp_con1}) and then divide each side by $CapF_{[1]}$, 
we get:
\begin{eqnarray}
\label{nonhomo-sum2}
 T_{[1]}+\frac{CapF_{[2]}}{CapF_{[1]}}T_{[2]}+..+ \frac{CapF_{[M]}}{CapF_{[1]}}T_{[M]} \geq \frac {\sum_{i \in N}d_i}{CapF_{[1]}} 
\end{eqnarray} 
If we let 
\begin{eqnarray}
 nREQ2 =\left \lceil \frac{\sum_{i \in N }d_i}{CapF_{[1]}} \right \rceil \nonumber 
\end{eqnarray}   
then, by rounding up the coefficients of inequality (\ref{nonhomo-sum2}), we get:
\begin{eqnarray}
\label{lb-open-fac2}
 \sum_{j \in M}T_j \geq nREQ2 
\end{eqnarray}  

The value of $nREQ_2$ is calculated based on the assumption that all facilities have the same capacity, which is 
the largest capacity value. However, in the calculation procedure of $nREQ$ for problems 
that have different capacity values (case 2), it is considered that the second facility might not have a capacity
as large as the first facility, the third facility might not have a capacity as large as the second facility and
so on.  Therefore, $nREQ \geq nREQ2$ and the constraint (\ref{lb-open-fac}) dominates constraint 
(\ref{lb-open-fac2}). If $\left(\frac{\sum_{i \in N }d_i}{CapF_{[1]}}\right)$ is fractional or 
$nREQ2<nREQ$, then the CG rank of constraint (\ref{lb-open-fac}) is 1. \qed
\end{proof}



Due to the similarity in structure between the SPP-LRS formulation and the CFLP, we might consider adding the set of constraints 
\begin{eqnarray}
Z_{jp} \leq & T_{j}  & \forall j \in M, \forall p \in P_{j} \label{relation-Z-T}
\end{eqnarray}
In the CFLP, constraints similar to inequality (\ref{relation-Z-T}) are facets of the CFL polytope if the
demands of the customers in the constraints are less than the capacity of the facilities (\citet{Leung}). 
In addition, \citet{Daskin} remarks that constraints in the CFLP similar to (\ref{relation-Z-T}) 
considerably improve the LP bound. 

Constraints (\ref{relation-Z-T})  ensure that a pairing of a facility can be selected only if that facility is 
selected. Although constraints (\ref{relation-Z-T}) are implied by constraints (\ref{spp_con2}),  
constraints (\ref{relation-Z-T}) may improve the LP relaxation of the model as in the case of 
CFLP. However, since there is potentially an exponential number of such constraints, we would
like to find a smaller set of constraints to add. 

\begin {prop} 
  To strengthen the LP relaxation bound, we can replace constraints (\ref{relation-Z-T}) with 
  \begin{eqnarray}
    \label{relation}
    \sum_{p \in P_{j}}a_{ipj}Z_{jp}\ \leq \ T_{j} \mbox{\ \ $\forall j \in M$ and $i \in N$} 
  \end{eqnarray}
and constraints (\ref{relation}) have a CG rank of 1.
\end{prop}
 
\begin{proof} 
Constraints (\ref{relation}) are valid for the integer program. In any feasible solution,
each demand node $i \in N$ is assigned to exactly one facility, so $\sum_{p \in P_j} a_{ipj} Z_{jp}$
is $1$ for one facility and $0$ for all others. For any $j \in M$ such that 
$\sum_{p \in P_j} a_{ipj} Z_{jp}$ is $1$, $T_j$ will be $1$, which corresponds 
to pairings being selected only for open facilities. 

%In any feasible integer solution, a demand node 
%$i \in N$ can be assigned to only one facility, so 
%$\sum_{p \in P_j} a_{ipj} Z_{jp}$ can only be 1 or 0 for any facility $j \in M$. Also 
%constraints (\ref{spp_con2}) guarantee that if any demand node is assigned to facility j,
%then $T_j$ must be greater than zero. Since $T_j$ is binary variable, if $T_j>0$, $T_j$ must be 1.
%Therefore, the inequality (\ref{relation}) is valid for any feasible integer solution. 

Also, when the coefficients of the variables in inequality (\ref{relation}) and constraint 
(\ref{relation-Z-T}) are compared, it is easy to see that inequality (\ref{relation}) 
dominates constraint (\ref{relation-Z-T}). For example, if we take $Z_{j'p'} \leq T_{j'}$
from constraint set (\ref{relation-Z-T}) for $j=j'$ and $p=p'$, we can compare it with a
constraint in set (\ref{relation}) for which $j=j'$ and $i=i'$ such that $a_{i'p'j'}=1$. 
That is $Z_{j'p'}+\sum_{p \in P_{j'}\setminus\{p'\}}a_{i'pj'}Z_{j'p}\ \leq \ T_{j'}$. Because 
of the $Z_{jp}$ terms other than  $Z_{j'p'}$ on the left side of the constraint,   
$\sum_{p \in P_{j'}}a_{i'pj'}Z_{j'p}\ \leq \ T_{j'}$ dominates $Z_{j'p'} \leq T_{j'}$. 

Now we show the rank of the inequality $\sum_{p \in P_j} a_{ipj} Z_{jp} \leq T_j$ 
for an arbitrary i and j. For notational convenience, suppose $i=1$ and $j=1$.  

Let $D_k=\sum_{i \in N}a_{ikj}d_i$ be the total demand of a pairing k . 
From constraint sets (\ref{spp_con1}) ($\leq$ part of the constraints) and (\ref{spp_con2}), 
we can write:
\begin{eqnarray}
 \sum_{\{p \in P_1 | a_{1p1}=1\}} Z_{1p} + \sum_{\{j \in M \setminus \{1\}\}}\sum_{\{p \in P_{j}|a_{1pj}=1\}}Z_{jp}  &\leq & 1 \label{prop2-a} \\
 \sum_{\{p \in P_1 | a_{1p1}=1\}} D_p Z_{1p} + \sum_{\{p \in P_1 | a_{1p1}=0\}} D_p Z_{1p} - CapF_1 T_1 & \leq & 0 \label{prop2-b} 
\end{eqnarray} 

Let $min_{\{p \in P_1 | a_{1p1}=1\}} \{ D_p \}=D_{min}$. Multiply inequality (\ref{prop2-a}) 
with $(CapF_1-D_{min})$ and add the resulting inequality to the inequality (\ref{prop2-b}). 
Then after dividing the obtained inequality by $CapF_1$ round it to floor: 

\begin{eqnarray}
 \sum_{\{p \in P_1 | a_{1p1}=1\}}  \left \lfloor \frac{CapF_1-D_{min}+D_p}{CapF_1} \right \rfloor Z_{1p} &+&
    \sum_{\{j \in M / \{1\}\}}\sum_{\{p \in P_{j}|a_{1pj}=1\}}\left \lfloor \frac{CapF_1-D_{min}}{CapF_1}\right \rfloor Z_{jp}\nonumber \\
     &+&\sum_{\{p \in P_1 | a_{1p1}=0\}}\left \lfloor \frac{D_p}{CapF_1}\right \rfloor Z_{1p} - T_1\nonumber \\
&\leq&  \left \lfloor \frac {CapF_1-D_{min}}{CapF_1}\right \rfloor  \label{prop2-c}
\end{eqnarray} 


Also, we simply assume that $D_p \leq CapF_1 \ \ \forall p \in P_1$ since otherwise we can set
$Z_{jp}=0$ and remove the variable from the constraint. This is true, because if the
total demand of the pairing is greater than the capacity of the facility, this pairing
cannot be chosen in a feasible solution for this facility. In inequality (\ref{prop2-c}), 
the coefficients in the first summation are in the range $[1,2)$, the coefficients in the second 
summation are in $[0,1)$ and if $D_p<CapF_1$ for all p in the third summation, the coefficients are in 
$[0,1)$. The right hand side is in the range $[0,1)$. Thus when we round down the coefficients, we get 
$\sum_{p \in P_1} a_{1p1} Z_{1p} \leq T_1$. Since we have obtained inequality (\ref{relation}) by using
the original constraints and the CG inequality procedure, the CG rank of the inequality is 1. Also, 
this inequality is not dominated by any linear combination of the original constraints, thus
its CG rank cannot be 0. If $D_p=CapF_1$ for any p in the third summation, when we round down, 
the coefficient for this variable in (\ref{prop2-c}) is 1. Thus, we must add the variables that 
satisfy $D_p=CapF_1$ to our inequality (\ref{relation}) to get a stronger constraint. But if we do not add these 
variables to the constraint (\ref{relation}), the constraint (\ref{relation}) will still be valid. 
\qed
\end{proof}

In chapter \ref{computationalexp}, we empirically demonstrate the improvement in the LP relaxation 
bounds achieved by constraints (\ref{lb-open-fac}), (\ref{relation-Z-T}) and (\ref{relation}). 
Throughout the remainder of the proposal, we use the augmented SPP-LRS formulation (AUG-SPP-LRS),
which includes the objective function (\ref{spp-obj}) and the constraints (\ref{spp_con1}), 
(\ref{spp_con2}), (\ref{binary1}), (\ref{binary2}), (\ref{lb-open-fac}), and (\ref{relation}). 


